% =====================================================================
%  Преамбула отчета о НИР (ГОСТ 7.32-2017)
%  Сборка: XeLaTeX + Biber
%  Стиль библиографии: ГОСТ Р 7.0.100-2018 (biblatex-gost)
% =====================================================================

% Кодировка и шрифты
\usepackage{fontspec}
\setmainfont{Times New Roman}
\setsansfont{Arial}
\setmonofont{Consolas}[Scale=0.92]
\newfontfamily\cyrillicfonttt{Consolas}[Scale=0.92]

\usepackage{polyglossia}
\setdefaultlanguage{russian}
\setotherlanguage{english}

% Геометрия страницы (ГОСТ 7.32-2017, п. 6.1):
% левое 30 мм, правое 15 мм, верхнее 20 мм, нижнее 20 мм
\usepackage[a4paper, left=30mm, right=15mm, top=20mm, bottom=20mm,
            footskip=10mm]{geometry}

% Межстрочный интервал 1.5 (ГОСТ 7.32-2017, п. 6.1)
\usepackage{setspace}
\onehalfspacing

% Абзацный отступ 1.25 см
\usepackage{indentfirst}
\setlength{\parindent}{1.25cm}

% Математика
\usepackage{amsmath, amssymb, amsthm, mathtools}

% Графика, таблицы, списки
\usepackage{graphicx}
\usepackage{booktabs}
\usepackage{longtable}
\usepackage{tabularx}
\usepackage{array}
\usepackage{multirow}
\usepackage{makecell}
\usepackage{caption}
\usepackage{subcaption}
\usepackage{float}
\usepackage{enumitem}

% Перечисления по ГОСТ 7.32-2017, п. 6.4.6:
%   itemize  --- маркер «тире» (—);
%   enumerate --- «1)», «2)», ... с абзацного отступа.
\setlist[itemize,1]{label=\textendash, leftmargin=1.25cm, labelsep=0.5em,
                     itemsep=0pt, parsep=0pt, topsep=2pt}
\setlist[itemize,2]{label=\textendash, leftmargin=2.5cm,  labelsep=0.5em,
                     itemsep=0pt, parsep=0pt, topsep=2pt}
\setlist[enumerate,1]{label=\arabic*), leftmargin=1.25cm, labelsep=0.5em,
                     itemsep=0pt, parsep=0pt, topsep=2pt}
\setlist[enumerate,2]{label=\arabic*), leftmargin=2.5cm,  labelsep=0.5em,
                     itemsep=0pt, parsep=0pt, topsep=2pt}

% Цвет, ссылки, листинги
\usepackage{xcolor}
\usepackage[hidelinks]{hyperref}
\hypersetup{
    pdfencoding=auto,
    bookmarksnumbered=true,
    bookmarksopen=true,
    colorlinks=false
}
\usepackage{listings}
\lstset{
    basicstyle=\ttfamily\footnotesize,
    breaklines=true,
    showstringspaces=false,
    keywordstyle=\bfseries,
    columns=fullflexible,
    captionpos=b,
    frame=single,
    aboveskip=8pt,
    belowskip=8pt
}

% Заголовки и оглавление
\usepackage{titlesec}
\usepackage{tocloft}

% Библиография по ГОСТ Р 7.0.100-2018
\usepackage[
    backend=biber,
    style=gost-numeric,
    language=auto,
    sorting=none,
    maxbibnames=4,
    minbibnames=1
]{biblatex}

% Дополнительные служебные пакеты
\usepackage{csquotes}
\usepackage{etoolbox}
\usepackage{url}
\usepackage[protrusion=true,expansion=false,
            final,kerning=false,spacing=false,tracking=false]{microtype}
\usepackage{xurl}
\usepackage{seqsplit}
\PassOptionsToPackage{hyphens}{url}
\sloppy
\emergencystretch=3em
\hbadness=10000
\vbadness=10000
\hfuzz=2pt
\vfuzz=2pt
\hyphenpenalty=1000
\tolerance=2000

% Сброс счетчиков таблиц, рисунков и формул при начале новой секции
\makeatletter
\@addtoreset{table}{section}
\@addtoreset{figure}{section}
\@addtoreset{equation}{section}
\makeatother

% =====================================================================
%  Отказ от курсива (ГОСТ 7.32-2017 строго ограничивает выделение)
% =====================================================================
% \emph и \textit преобразуются в обычный текст. \eng - то же самое.
\let\origemph\emph
\renewcommand{\emph}[1]{#1}
\let\origtextit\textit
\renewcommand{\textit}[1]{#1}

% =====================================================================
% Заголовки разделов и подразделов (ГОСТ 7.32-2017, п. 6.2.3):
% с абзацного отступа, с прописной буквы, полужирным шрифтом,
% не подчеркивая, без точки в конце.
% =====================================================================
\titleformat{\section}
    {\normalfont\fontsize{14}{17}\bfseries\centering}
    {\thesection}{1em}{}
\titlespacing*{\section}{\parindent}{18pt}{12pt}

\titleformat{\subsection}
    {\normalfont\fontsize{14}{17}\bfseries}
    {\thesubsection}{1em}{}
\titlespacing*{\subsection}{\parindent}{14pt}{10pt}

\titleformat{\subsubsection}
    {\normalfont\fontsize{14}{17}\bfseries}
    {\thesubsubsection}{1em}{}
\titlespacing*{\subsubsection}{\parindent}{12pt}{8pt}

% Нумерация рисунков, таблиц, формул --- в пределах раздела
\renewcommand{\thefigure}{\arabic{section}.\arabic{figure}}
\renewcommand{\thetable}{\arabic{section}.\arabic{table}}
\renewcommand{\theequation}{\arabic{section}.\arabic{equation}}

% =====================================================================
% Подписи рисунков и таблиц (ГОСТ 7.32-2017, п. 6.5.7 и 6.6.3)
% =====================================================================
\DeclareCaptionLabelSeparator{gostdash}{~---~}

% «Рисунок N --- Наименование» по центру под рисунком, без точки.
\captionsetup[figure]{
    labelsep=gostdash, name=Рисунок,
    justification=centering, singlelinecheck=true,
    font=normalsize, position=below
}
% «Таблица N --- Наименование» слева над таблицей, без абзацного отступа,
% без точки в конце.
\captionsetup[table]{
    labelsep=gostdash, name=Таблица,
    justification=raggedright, singlelinecheck=false,
    font=normalsize, position=above, skip=4pt
}
\captionsetup[longtable]{
    labelsep=gostdash, name=Таблица,
    justification=raggedright, singlelinecheck=false,
    font=normalsize, position=above, skip=4pt
}

% =====================================================================
% Стиль таблиц (ГОСТ 7.32-2017, п. 6.6.6):
% «Таблицы слева, справа, сверху и снизу ограничивают линиями».
% Заменяем booktabs-команды \toprule/\midrule/\bottomrule на \hline,
% вертикальные линии указываются в спецификации таблицы ( |...| ).
% Заголовки граф выравниваются по центру (п. 6.6.6), данные --- по левому
% краю или по центру в зависимости от содержимого.
% =====================================================================
\renewcommand{\toprule}{\hline}
\renewcommand{\midrule}{\hline}
\renewcommand{\bottomrule}{\hline}
\renewcommand{\arraystretch}{1.15}
\setlength{\tabcolsep}{6pt}

% Колонки с переносом по словам: L --- по левому краю; C --- центр.
\newcolumntype{L}[1]{>{\raggedright\arraybackslash}p{#1}}
\newcolumntype{C}[1]{>{\centering\arraybackslash}p{#1}}
\newcolumntype{Y}{>{\centering\arraybackslash}X}
\newcolumntype{Z}{>{\raggedright\arraybackslash}X}

% Колонтитулы
\usepackage{fancyhdr}
\pagestyle{fancy}
\fancyhf{}
\fancyfoot[C]{\thepage}
\renewcommand{\headrulewidth}{0pt}
\renewcommand{\footrulewidth}{0pt}

% =====================================================================
% СОДЕРЖАНИЕ (ГОСТ 7.32-2017, п. 5.4, 6.13):
% Каждую запись --- отдельным абзацем, выровненным влево; номера страниц
% выровнены по правому краю, соединены отточием. Обозначения
% подразделов --- с отступом, равным двум знакам относительно
% обозначения разделов; пунктов --- четырем знакам.
% =====================================================================
\setcounter{tocdepth}{3}

% Сам заголовок «СОДЕРЖАНИЕ» --- центрированный, полужирный, прописными
\AtBeginDocument{%
  \renewcommand{\contentsname}{\hfill\bfseries СОДЕРЖАНИЕ\hfill\mbox{}}%
}

% Шрифт записей в TOC --- нормальный (по ГОСТу не выделяется)
\renewcommand{\cftsecfont}{\normalfont}
\renewcommand{\cftsecpagefont}{\normalfont}
\renewcommand{\cftsubsecfont}{\normalfont}
\renewcommand{\cftsubsecpagefont}{\normalfont}
\renewcommand{\cftsubsubsecfont}{\normalfont}
\renewcommand{\cftsubsubsecpagefont}{\normalfont}

% Отточия между текстом и номером страницы
\renewcommand{\cftdotsep}{1.5}
\renewcommand{\cftsecdotsep}{1.5}
\renewcommand{\cftsecleader}{\cftdotfill{\cftsecdotsep}}

% Отступы: подраздел --- начинается с того же положения, что и текст
% названия раздела (т.е. сдвинут на ширину номера раздела); пункт ---
% сдвинут на ширину номера подраздела.
\setlength{\cftsecindent}{0pt}
\setlength{\cftsecnumwidth}{1.6em}
\setlength{\cftsubsecindent}{1.6em}
\setlength{\cftsubsecnumwidth}{2.4em}
\setlength{\cftsubsubsecindent}{4.0em}
\setlength{\cftsubsubsecnumwidth}{3.2em}

% =====================================================================
% Пользовательские команды
% =====================================================================
\newcommand{\figref}[1]{рисунок~\ref{#1}}
\newcommand{\tabref}[1]{таблица~\ref{#1}}
\newcommand{\secref}[1]{раздел~\ref{#1}}
\newcommand{\eng}[1]{#1}

% Структурный элемент отчета по ГОСТ 7.32-2017 (п. 6.2.1):
% располагается в середине строки, без точки в конце, прописными буквами,
% не подчеркивая. Каждый структурный элемент --- с новой страницы.
\newcommand{\structelem}[1]{%
  \phantomsection%
  \addcontentsline{toc}{section}{#1}%
  \begin{center}\bfseries #1\end{center}%
  \vspace{6pt}%
}

% Структурный элемент БЕЗ записи в Содержание (для реферата,
% поскольку п. 5.4.1 ГОСТ 7.32-2017 не относит реферат к
% включаемым в содержание элементам).
\newcommand{\structelemnotoc}[1]{%
  \begin{center}\bfseries #1\end{center}%
  \vspace{6pt}%
}

\DeclareMathOperator*{\argmax}{arg\,max}
\DeclareMathOperator*{\argmin}{arg\,min}
\DeclareMathOperator{\softmax}{softmax}
\DeclareMathOperator{\KL}{KL}
\DeclareMathOperator{\MMD}{MMD}
\DeclareMathOperator{\PSI}{PSI}

% Метка последней страницы для реферата
\usepackage{lastpage}

% Заголовок списка источников --- структурный элемент.
\defbibheading{structelem}[\bibname]{%
  \phantomsection%
  \addcontentsline{toc}{section}{#1}%
  \begin{center}\bfseries #1\end{center}%
  \vspace{6pt}%
}
